% abstract.tex
\placeholder{Abstract -- rewrite in finished prose once Phase 2 sections land.}

\noindent
The A=1 Discrete Causal Lattice (DCL) framework establishes the
bipartite octahedral lattice as the substrate of fundamental physics
in (3+1) dimensions (\emph{Geometry First}, Paper~I).  Downstream
work in the DCL series -- the gauge-derivation programme of Paper~II,
the discrete-probability paper closing the gauge-coupling prefactor,
the planned Hilbert-Sixth axiomatisation capstone, and future
kinetic-theory work in the spirit of the Boltzmann / Navier-Stokes
limit -- consumes \emph{more} of the underlying mathematics than
Paper~I exposed, and uses it in arbitrary dimension rather than only
in the (3+1)-D physical setting.  This paper collects the
foundational mathematics those downstream subprojects share, in a
single citable source.

\medskip

\noindent
We define the bipartite octahedral lattice in arbitrary dimension
$n$, characterise its automorphism algebra as a function of $n$
(extending the 3-D 71-dimensional per-site algebra catalogued in
\texttt{dcl-zoo}), develop the discrete differential geometry,
function-space theory, and discrete measure theory the lattice
supports, and -- centrally -- formulate the discrete-to-continuum
limit operators that the Hilbert-Sixth capstone consumes for its
axiomatisation claim.  Two motivating examples appear in
Section~\ref{sec:motivating_examples}: Boltzmann-shaped collision
operators on the lattice, and the lattice-Boltzmann -> Navier-Stokes /
Euler fluid-limit ansatz.  In this first paper these examples are
framing only; full proofs are deferred to subsequent papers in this
repository or in sibling repositories.

\medskip

\noindent
The paper's status across its claims is tracked in
Section~\ref{sec:introduction}, Table~\ref{tab:audit}.  At v0.1-DRAFT,
all toolkit-internal rows are \texttt{STUB}; the artefact at this
version is a paper-shaped outline that locks the conventions and
the sectioning, not a finished body of proofs.
